\section{Description entropy in arbitrary matrix dimension}
\label{sec:matrixentropy76}
The intrinsic comparison of Theorem~\ref{thm:matrixmetric76} permits
the full effect body to be treated without choosing eigenvector
coordinates near repeated spectra. Throughout this section the input
dimension $d\ge1$ is fixed. Constants may depend on $d$, but not on the
integer horizon $N\ge1$ or on the accuracy parameter.

Write $\mathcal C_N(\mathfrak E_d,\delta)$ for the covering number under
$d_N$, with centres in the same legal memoryless ordered binary
interface, and write $\mathcal C_N^{\mathrm{rat}}$ when the real and
imaginary parts of every matrix entry of each centre are rational.

\begin{theorem}[Entropy of the full effect body]\label{thm:matrixcover76}
For every integer $d\ge1$ there are constants $c_d,C_d,\delta_d>0$
such that, for all integers $N\ge1$ and all $0<\delta\le\delta_d$,
\begin{equation}\label{eq:matrixcover76}
\begin{split}
 c_d N^{d^2/2}[\log(N+2)]^{\lfloor d/2\rfloor}\delta^{-d^2}
 &\le \mathcal C_N(\mathfrak E_d,\delta)\\
 &\le \mathcal C_N^{\mathrm{rat}}(\mathfrak E_d,\delta)\\
 &\le C_d N^{d^2/2}[\log(N+2)]^{\lfloor d/2\rfloor}\delta^{-d^2}.
\end{split}
\end{equation}
Consequently the optimum fixed-length reusable description has length
\begin{equation}\label{eq:matrixbits76}
 \frac{d^2}{2}\log_2N
 +\lfloor d/2\rfloor\log_2\log(N+2)
 +d^2\log_2(1/\delta)+O_d(1).
\end{equation}
The estimates include all support and spectral multiplicities. The
rational-centre assertion is an existence statement in every dimension;
Theorem~\ref{thm:biasedcodec75} gives the separate exact encoder in
dimension two.
\end{theorem}

For $d=1$, the result is the Bernoulli covering law
$N^{1/2}\delta^{-1}$. For $d=2$, it recovers
Theorem~\ref{thm:biasedcover75}. The next two cases have orders
$N^{9/2}\log(N+2)\delta^{-9}$ and
$N^8[\log(N+2)]^2\delta^{-16}$. The logarithmic power counts independent
nested approaches of paired eigenvalues to the two support endpoints.
We first establish a uniform measure estimate for actual operational
balls, and then evaluate the total measure.

\subsection{Operational balls and a regularized volume form}
Set $n=d^2$, $\tau=N^{-1}$ and
\[
 W_E=E(I-E)+\frac\tau2 I.
\]
As in Section~\ref{sec:matrixmetric76}, let
\begin{equation}\label{eq:matrixfixedform76}
 g_{N,E}(H)^2
 =N\langle H,(L_{W_E}+R_{W_E})^{-1}H\rangle_{\mathrm{HS}}
\end{equation}
on the real Hilbert space of Hermitian matrices. Let $G_{N,E}$ denote
this positive quadratic form in Hilbert--Schmidt orthonormal
coordinates, and define
\begin{equation}\label{eq:matrixmeasure76}
 \rho_N(E)=\sqrt{\det G_{N,E}},\qquad
 d\mu_N(E)=\rho_N(E)\,dE,
\end{equation}
where $dE$ is Lebesgue measure in those coordinates. The cutoff makes
$\rho_N$ positive and continuous on the whole closed body for each
$N$.

\begin{lemma}[Uniform volume of operational balls]\label{lem:matrixball76}
For fixed $d$ there are $a_d,A_d,r_d>0$ such that
\begin{equation}\label{eq:matrixball76}
 a_dr^{d^2}\le
 \mu_N\bigl(B_{d_N}(E,r)\bigr)
 \le A_dr^{d^2}
\end{equation}
for every $N\ge1$, every $E\in\mathfrak E_d$, and every
$0<r\le r_d$. Balls are taken inside $\mathfrak E_d$.
\end{lemma}
\begin{proof}
We make explicit the two local consequences of
Lemma~\ref{lem:matrixlocal76} that will be used. If $H$ is Hermitian and
$q=g_{N,E}(H)$, diagonalizing $W_E$ gives
\begin{equation}\label{eq:matrixnormalized76}
 \|W_E^{-1/2}HW_E^{-1/2}\|_{\mathrm{HS}}^2
 \le4q^2,
\end{equation}
since its eigenvalues $w_i$ are at least $\tau/2$ and
\[
 \frac{w_i+w_j}{Nw_iw_j}
                  =\tau(w_i^{-1}+w_j^{-1})\le4.
\]
For any legal $E+tH$, direct expansion yields
\[
 W_{E+tH}-W_E
 =\frac t2\bigl((I-2E)H+H(I-2E)\bigr)-t^2H^2.
\]
The matrices $E$ and $W_E$ commute, $\|I-2E\|_{\rm op}\le1$, and
$\|W_E\|_{\rm op}\le3/4$. Thus the norm of the normalized difference
is at most $2|t|q+3t^2q^2$. For $|t|\le1$ and $q$ below an absolute
constant, it follows that
\begin{equation}\label{eq:matrixvariancecomparison76}
 \tfrac12W_E\preceq W_{E+tH}\preceq2W_E.
\end{equation}
Loewner comparison passes to $L_W+R_W$ and, by inversion, to the
quadratic forms. It also gives
\begin{equation}\label{eq:matrixdensitycomparison76}
 2^{-n/2}\rho_N(E)\le\rho_N(E+tH)\le2^{n/2}\rho_N(E).
\end{equation}

Apply this comparison from an endpoint to the midpoint and from the
midpoint to either endpoint. Theorem~\ref{thm:matrixmetric76} then
provides constants $A,B,r_0>0$, depending only on $d$, such that
\begin{align}
 d_N(E,F)\le r&\ \Longrightarrow\
       g_{N,E}(F-E)\le Ar,\label{eq:matrixballupperellipse76}\\
 g_{N,E}(F-E)\le Br&\ \Longrightarrow\
       d_N(E,F)\le r,\label{eq:matrixballlowerellipse76}
\end{align}
for legal $E,F$ and $0<r\le r_0$. Decrease $r_0$ so that $Ar_0$
lies in the small-form range of \eqref{eq:matrixvariancecomparison76}.

Let $\omega_n$ be the Euclidean volume of the unit ball in $\R^n$.
The ellipsoid $g_{N,E}(H)\le u$ has Lebesgue volume
$\omega_nu^n/\rho_N(E)$. Equations
\eqref{eq:matrixballupperellipse76} and
\eqref{eq:matrixdensitycomparison76} therefore imply
\[
 \mu_N\bigl(B_{d_N}(E,r)\bigr)
                  \le2^{n/2}\omega_n(Ar)^n.
\]

For the lower bound we construct a full legal ellipsoid, including
when $E$ lies on a support boundary. Choose a small constant $a>0$
depending only on $d$, set $s=ar\tau$, and move towards the scalar
interior point:
\begin{equation}\label{eq:matrixinwardshift76}
 E_s=(1-2s)E+sI.
\end{equation}
Take $r_0$ sufficiently small that $s\le1/4$. Then
$sI\preceq E_s\preceq(1-s)I$. Since the displacement commutes with
$E$, its form is bounded by
\[
 g_{N,E}(E_s-E)^2
 =N\sum_i\frac{s^2(1-2\lambda_i)^2}
                    {2\lambda_i(1-\lambda_i)+\tau}
 \le d\frac{s^2}{\tau^2}=da^2r^2.
\]
Choose $a\sqrt d\le B/4$. By
\eqref{eq:matrixballlowerellipse76},
$d_N(E,E_s)\le r/4$.

Let $X=X^*$ satisfy $g_{N,E_s}(X)\le br$, with $b>0$ to be chosen.
Write $\ell_i$ for the eigenvalues of $E_s$ and
$w_i=\ell_i(1-\ell_i)+\tau/2$. Since $\ell_i\ge s$ and
$w_i\le\ell_i+\tau/2$, we have
\begin{align*}
 \|E_s^{-1/2}XE_s^{-1/2}\|_{\mathrm{HS}}^2
 &\le g_{N,E_s}(X)^2
       \max_{i,j}\frac{w_i+w_j}{N\ell_i\ell_j}\\
 &\le b^2r^2\left(\frac{2\tau}{s}
                         +\frac{\tau^2}{s^2}\right)
   =b^2\left(\frac{2r}{a}+\frac1{a^2}\right).
\end{align*}
The same bound holds with $I-E_s$ in place of $E_s$, because
$1-\ell_i\ge s$ and $w_i\le1-\ell_i+\tau/2$. Taking
$b\le a/4$ and $ar_0\le1$ makes both operator norms at most $1/2$.
Consequently every point of the entire ellipsoid
\[
 \mathcal A_s=\{E_s+X:g_{N,E_s}(X)\le br\}
\]
is a legal effect. Choose also $b\le B/4$. Equation
\eqref{eq:matrixballlowerellipse76} and the triangle inequality give
$\mathcal A_s\subset B_{d_N}(E,r)$. The density comparison about
$E_s$ and the exact ellipsoid volume now imply
\[
 \mu_N\bigl(B_{d_N}(E,r)\bigr)
           \ge2^{-n/2}\omega_n(br)^n.
\]
All choices depend only on $d$. This proves the assertion on every
boundary stratum as well as in the interior.
\end{proof}

\subsection{The total regularized volume}
\begin{proposition}[Nested endpoint accumulation]\label{prop:matrixvolume76}
For fixed $d\ge1$, uniformly over $N\ge1$,
\begin{equation}\label{eq:matrixvolume76}
 \mu_N(\mathfrak E_d)
       \asymp_d N^{d^2/2}[\log(N+2)]^{\lfloor d/2\rfloor}.
\end{equation}
\end{proposition}
\begin{proof}
In a spectral basis, the quadratic form has weight $N/(2w_i)$ on
the $i$th diagonal direction and weight $N/(w_i+w_j)$ on each of
the two Hilbert--Schmidt orthonormal real directions belonging to
$i<j$. Hence
\begin{equation}\label{eq:matrixdeterminant76}
 \rho_N(E)=2^{-d/2}N^{d^2/2}
       \prod_iw_i^{-1/2}\prod_{i<j}(w_i+w_j)^{-1}.
\end{equation}
We use the classical spectral change of variables
\cite{SZGeom76}. For a simple spectrum, differentiation of
$E=U\operatorname{diag}(\lambda)U^*$ gives
\[
 U^*(dE)U=d\operatorname{diag}(\lambda)
       +[U^*dU,\operatorname{diag}(\lambda)].
\]
Each complex off-diagonal pair therefore contributes
$(\lambda_i-\lambda_j)^2$ to the real Jacobian. Integration over
the compact unitary orbit and over permutations supplies a positive
constant depending only on $d$. The repeated-spectrum set is
Lebesgue null. This spectral change of variables, together with
\eqref{eq:matrixdeterminant76}, yields
\begin{equation}\label{eq:matrixspectralintegral76}
 \mu_N(\mathfrak E_d)\asymp_d N^{d^2/2}I_d(\tau),
\end{equation}
where
\begin{equation}\label{eq:matrixId76}
 I_d(\tau)=\int_{[0,1]^d}
 \prod_i(v_i+\tau)^{-1/2}
 \prod_{i<j}\frac{(\lambda_i-\lambda_j)^2}{v_i+v_j+\tau}
 \,d\lambda,\qquad v_i=\lambda_i(1-\lambda_i).
\end{equation}
Replacing $v_i+\tau/2$ by $v_i+\tau$ in the single-coordinate
factors changes the integral by a factor between fixed positive
constants depending only on $d$.

We prove
\begin{equation}\label{eq:matrixintegralasymptotic76}
 I_d(\tau)\asymp_d
       [\log(2+\tau^{-1})]^{\lfloor d/2\rfloor},
                       \qquad0<\tau\le1.
\end{equation}
For the upper bound, partition each coordinate into
$[0,1/4]$, $[1/4,3/4]$, and $[3/4,1]$.
Factors involving an eigenvalue in the middle interval are bounded
above by constants, so these coordinates can be integrated out.
For each of the remaining $l\le d$ coordinates put $x_i=\lambda_i$
near zero and $x_i=1-\lambda_i$ near one. On a sector
$0\le x_1\le\cdots\le x_l\le1/4$, let
$\epsilon_i\in\{1,-1\}$ record the endpoint and put
\[
 S_j=\sum_{i=1}^j\epsilon_i,\qquad S_0=0,\qquad y_j=x_j+\tau.
\]
The individual factor is at most $Cy_j^{-1/2}$. For $i<j$, a
same-endpoint pair factor is at most $Cy_j$, whereas an
opposite-endpoint factor is at most $C/y_j$. The sector integrand
is consequently bounded by
\begin{equation}\label{eq:matrixsector76}
 C_d\prod_{j=1}^ly_j^{a_j-1},\qquad
 a_j=\frac12+\epsilon_jS_{j-1}
              =\frac{S_j^2-S_{j-1}^2}{2}.
\end{equation}
Its cumulative exponents satisfy
\begin{equation}\label{eq:matrixprefix76}
 A_j=\sum_{i=1}^ja_i=\frac{S_j^2}{2}\ge0.
\end{equation}

Set $T=1/4+\tau$ and $L=\log(T/\tau)$. Substitute
$y_j=T\exp(-z_j)$, where $z_1\ge\cdots\ge z_l\ge0$,
and then $u_j=z_j-z_{j+1}$ for $j<l$, $u_l=z_l$.
The sector integral is bounded by
\begin{equation}\label{eq:matrixprefixintegral76}
 C_d\int_{u_j\ge0,\,\sum_j u_j\le L}
              \exp\left(-\sum_jA_ju_j\right)du.
\end{equation}
The omitted factor $T^{A_l}$ is bounded above uniformly in
$0<\tau\le1$. Each $A_j>0$ is at least $1/2$, so its variable has
an integrable exponential tail. Each vanishing $A_j$ contributes
at most one factor $1+L$. A vanishing prefix sum requires an even
$j$, and there are at most $\lfloor l/2\rfloor$ such indices.
Summing the finitely many sectors proves the upper estimate in
\eqref{eq:matrixintegralasymptotic76}. The case $l=0$ contributes
only a bounded constant.

For the lower bound let $m=\lfloor d/2\rfloor$. Choose scales
\begin{equation}\label{eq:matrixnestedscales76}
 \tau\le t_1,\qquad 8t_j\le t_{j+1}\ (j<m),
                 \qquad t_m\le1/32,
\end{equation}
and restrict $2m$ eigenvalues in pairs by
\begin{equation}\label{eq:matrixpairedboxes76}
 \lambda_j^-=a_j\in[t_j,2t_j],\qquad
 \lambda_j^+=1-b_j,\quad b_j\in[t_j,2t_j],
                         \quad1\le j\le m.
\end{equation}
If $d$ is odd, restrict the remaining eigenvalue to $[3/8,5/8]$;
all its factors are bounded above and below by positive constants.
For one pair the two individual spectral factors and its
opposite-endpoint factor have product comparable to $t_j^{-2}$.
For two pairs at scales $t_i\le t_j/8$, the two same-endpoint
factors are comparable to $t_j$ and the two opposite-endpoint
factors to $t_j^{-1}$. Their product is therefore bounded above
and below by fixed constants. Thus the integrand on the box is
comparable to $\prod_jt_j^{-2}$ and its volume to
$\prod_jt_j^2$. Each box contributes at least $c_d>0$.

Take the scales from $\tau,8\tau,8^2\tau,\ldots$ up to $1/32$.
The boxes for distinct increasing $m$-tuples are disjoint.
Their number is a binomial coefficient bounded below by
$c_d[\log(\tau^{-1})]^m$ once $\log(\tau^{-1})$ is sufficiently
large depending on $d$. In the remaining bounded range, a fixed
box of distinct eigenvalues strictly inside $(0,1)$ has a
positive uniform integral and supplies the lower estimate.
For $m=0$ this interior argument suffices on the full range.
This proves \eqref{eq:matrixintegralasymptotic76} and, by
\eqref{eq:matrixspectralintegral76}, the proposition.
\end{proof}

The balanced-prefix identity \eqref{eq:matrixprefix76} accounts for
all nested endpoint scales. A simultaneous approach of all eigenvalues
to the endpoints detects only one possible critical scaling; the
separated pairs in \eqref{eq:matrixpairedboxes76} show why the
logarithmic multiplicity can be larger.

\subsection{Covering, rational centres, and description length}
\begin{proof}[Proof of Theorem~\ref{thm:matrixcover76}]
Choose $\delta_d$ smaller than the radius in
Lemma~\ref{lem:matrixball76}. A cover by $L$ balls of radius
$\delta$ satisfies
\[
 \mu_N(\mathfrak E_d)\le LA_d\delta^{d^2},
\]
which gives the required lower bound. Conversely, take a maximal
$\delta$-separated subset of $\mathfrak E_d$. Its balls of radius
$\delta/3$ are pairwise disjoint, and the lower bound in
\eqref{eq:matrixball76} bounds its cardinality by
\[
 C_d\mu_N(\mathfrak E_d)\delta^{-d^2}.
\]
It is a radius-$\delta$ cover by maximality. Such finite maximal
sets exist: the effect body is compact, $d_N$ is continuous by the
hybrid bound, and $d_1(E,F)=2\|E-F\|_{\rm op}$ separates distinct
effects. Proposition~\ref{prop:matrixvolume76} now gives both
estimates for legal centres.

Rational legal effects are dense. Indeed, first move towards $I/2$
by \eqref{eq:matrixinwardshift76}, and then approximate all independent
real matrix coordinates by rational numbers while remaining strictly
between $0$ and $I$. For fixed $N$, continuity in the operational
distance follows again from the hybrid bound. Replace every centre
of a radius-$\delta/2$ cover by a rational legal effect at distance
less than $\delta/2$. The resulting radius-$\delta$ cover has the
same cardinality. This proves the rational-centre upper bound.

A fixed-length payload indexes the centres of a cover, and conversely
each legal reusable program with a fixed payload supplies a centre.
Taking base-two logarithms of \eqref{eq:matrixcover76} and rounding
the index length gives \eqref{eq:matrixbits76}.
\end{proof}
